\documentclass[tikz,border=8pt]{standalone}
\usepackage{tikz}
\usepackage{amsmath}
\usepackage{pgfplots}
\pgfplotsset{compat=1.18}
\usepackage{ctex}
\usetikzlibrary{calc, positioning, arrows.meta, shapes, fit, backgrounds}

\begin{document}
\begin{tikzpicture}

% === 左子图：均匀分布 H = log4 = 2 bit ===
\begin{axis}[
  name=ax1,
  width=5.2cm, height=4.5cm,
  xmin=0.3, xmax=4.7,
  ymin=0, ymax=0.5,
  xlabel={符号},
  ylabel={概率},
  axis lines=left,
  xtick={1,2,3,4},
  xticklabels={$x_1$,$x_2$,$x_3$,$x_4$},
  xticklabel style={font=\small},
  yticklabel style={font=\small},
  ytick={0,0.1,0.2,0.3,0.4},
  ymax=0.45,
  grid=major, grid style={gray!30},
  clip=false
]
  \addplot+[ybar, bar width=0.5, fill=blue!40, draw=blue!70!black, mark=none] coordinates {
    (1,0.25) (2,0.25) (3,0.25) (4,0.25)
  };
\end{axis}
\node[font=\small\bfseries] at ($(ax1.north)+(0,0.35cm)$) {均匀分布};
\node[font=\scriptsize, fill=white, draw=gray!50, thin, rounded corners=2pt, inner sep=3pt]
  at ($(ax1.south)+(0,-0.8cm)$) {$H=\log 4=2$ bit};

% === 中子图：双峰分布 H ≈ 1.5 bit ===
\begin{axis}[
  name=ax2,
  at={(ax1.east)}, anchor=west, xshift=0.8cm,
  width=5.2cm, height=4.5cm,
  xmin=0.3, xmax=4.7,
  ymin=0, ymax=0.55,
  xlabel={符号},
  ylabel={},
  axis lines=left,
  xtick={1,2,3,4},
  xticklabels={$x_1$,$x_2$,$x_3$,$x_4$},
  xticklabel style={font=\small},
  yticklabel style={font=\small},
  ytick={0,0.1,0.2,0.3,0.4,0.5},
  grid=major, grid style={gray!30},
  clip=false
]
  \addplot+[ybar, bar width=0.5, fill=orange!50, draw=orange!80!black, mark=none] coordinates {
    (1,0.45) (2,0.05) (3,0.45) (4,0.05)
  };
\end{axis}
\node[font=\small\bfseries] at ($(ax2.north)+(0,0.35cm)$) {双峰分布};
\node[font=\scriptsize, fill=white, draw=gray!50, thin, rounded corners=2pt, inner sep=3pt]
  at ($(ax2.south)+(0,-0.8cm)$) {$H\approx 1.39$ bit};

% === 右子图：单峰分布 H ≈ 0.8 bit ===
\begin{axis}[
  name=ax3,
  at={(ax2.east)}, anchor=west, xshift=0.8cm,
  width=5.2cm, height=4.5cm,
  xmin=0.3, xmax=4.7,
  ymin=0, ymax=0.85,
  xlabel={符号},
  ylabel={},
  axis lines=left,
  xtick={1,2,3,4},
  xticklabels={$x_1$,$x_2$,$x_3$,$x_4$},
  xticklabel style={font=\small},
  yticklabel style={font=\small},
  ytick={0,0.2,0.4,0.6,0.8},
  grid=major, grid style={gray!30},
  clip=false
]
  \addplot+[ybar, bar width=0.5, fill=green!45!black!60, draw=green!50!black, mark=none] coordinates {
    (1,0.7) (2,0.1) (3,0.1) (4,0.1)
  };
\end{axis}
\node[font=\small\bfseries] at ($(ax3.north)+(0,0.35cm)$) {单峰分布};
\node[font=\scriptsize, fill=white, draw=gray!50, thin, rounded corners=2pt, inner sep=3pt]
  at ($(ax3.south)+(0,-0.8cm)$) {$H\approx 0.80$ bit};

% 整体标题
\node[font=\large\bfseries] at ($(ax1.north)!0.5!(ax3.north)+(0,1.2cm)$)
  {熵 $H(X)$ 的几何示意：概率集中度越高，熵越小};

% 底部公式
\node[font=\normalsize, fill=white, draw=gray!50, thin, rounded corners=2pt, inner sep=5pt]
  at ($(ax1.south)!0.5!(ax3.south)-(0,1.6cm)$)
  {$H(X)=-\displaystyle\sum_{i} p_i \log p_i$};

\end{tikzpicture}
\end{document}
